This section defines a reproducible protocol for retrieving, ranking, and injecting evidence for each section: focused query formation, lightweight candidate scoring and ranking, conservative selection and prompt injection, and explicit provenance recording. Queries combine the section title and seed notes, local predecessor summaries, and task constraints; retrieval is scoped to a user-specified corpus and enforces minimum snippet lengths to reduce topical drift. Retrieved snippets are scored by a configurable similarity metric, filtered by simple acceptability checks (presence of a bibliographic identifier, ability to extract a quote or paraphrase, and absence of OCR errors), and collapsed into a compact, non-duplicative set that is injected into prompts with deterministic labels and provenance. Prompts include a JSON schema requiring fact text, supporting item id(s), and confidence; generators are instructed to cite supporting items and avoid inventing ungrounded claims, while the system emits minimal metadata for each injected item (source_id, source_type, page_or_location, snippet_text, context, extraction_method), applies validation and overlap-collapsing, falls back to evidence-missing placeholders when no acceptable candidates exist, and is modular, audit-friendly, and intentionally conservative (trading recall for precision at the cost of dependence on corpus coverage).